\documentclass[11pt,a4paper]{report}
% preamble.tex — Dossier scoring de confiance IDeXtend
% Style : report 11pt, Computer Modern, proche du rendu LaTeX classique

\usepackage[utf8]{inputenc}
\usepackage[T1]{fontenc}
\usepackage[french]{babel}
\usepackage{csquotes}

% Mise en page
\usepackage[a4paper, top=2.5cm, bottom=2.5cm, left=3cm, right=2.5cm]{geometry}

% Mathématiques
\usepackage{amsmath, amssymb}

% Tableaux
\usepackage{booktabs}
\usepackage{array}
\usepackage{longtable}

% Figures
\usepackage{graphicx}

% Couleurs et surlignage
\usepackage[table]{xcolor}
\usepackage{soul}

% Commandes de surlignage (compatibles avec soul)
\newcommand{\hlConnu}[1]{\sethlcolor{gray!40}\hl{#1}}
\newcommand{\hlNouveau}[1]{\sethlcolor{green!35}\hl{#1}}
\newcommand{\hlHorsSujet}[1]{\sethlcolor{yellow!60}\hl{#1}}
\newcommand{\hlInverifiable}[1]{\sethlcolor{orange!50}\hl{#1}}
\newcommand{\hlIncorrect}[1]{\sethlcolor{red!40}\hl{#1}}

% Code source
\usepackage{listings}
\lstset{
  basicstyle=\small\ttfamily,
  breaklines=true,
  frame=single,
  numbers=left,
  numberstyle=\tiny,
  backgroundcolor=\color{gray!8},
  xleftmargin=1.5em,
  framexleftmargin=1.5em,
}

% Hyperliens
\usepackage{hyperref}
\hypersetup{
  colorlinks=true,
  linkcolor=black,
  citecolor=black,
  urlcolor=black,
}

% En-têtes et pieds de page
\usepackage{fancyhdr}
\pagestyle{fancy}
\fancyhf{}
\fancyhead[LE,RO]{\thepage}
\fancyhead[RE]{\leftmark}
\fancyhead[LO]{\rightmark}
\renewcommand{\headrulewidth}{0.4pt}

% Encadrés simples (tcolorbox n'est pas requis — on utilise minipage + fbox)
\newenvironment{encadre}[1][]{%
  \medskip\noindent
  \begin{minipage}{\linewidth}
  \hrule\medskip
  \ifx&#1&\else{\small\textit{#1}}\par\smallskip\fi
}{%
  \medskip\hrule
  \end{minipage}\medskip
}


\title{Dossier technique\\[0.4em]
\large Scoring de confiance des composants IA\\[0.2em]
\large Système IDeXtend}
\author{Leveque Rayan}
\date{Mai 2026}

\begin{document}

% --- Page de garde --------------------------------------------------------
\begin{titlepage}
\centering
\vspace*{2cm}
{\Large Gendarmerie nationale \par}
\vspace{0.3cm}
{\large Centre Forensique d'Intelligence Artificielle (CFIA) \par}
\vspace{0.3cm}
{\large Leveque Rayan \par}
\vfill
{\Huge\bfseries Scoring de confiance\\[0.4em]
des composants IA\par}
\vspace{0.8cm}
{\Large Système IDeXtend\par}
\vfill
{\large\textit{Conformité AI Act --- Annexe III §6\\
Système à haut risque --- Application de la loi}\par}
\vspace{1cm}
{\Large\textbf{Version 0.1 --- Brouillon}\par}
\vspace*{2cm}
\end{titlepage}

\setcounter{tocdepth}{0}
\tableofcontents
\clearpage

% --- Introduction ---------------------------------------------------------
\chapter*{Introduction}
\addcontentsline{toc}{chapter}{Introduction}

IDeXtend est un système RAG (\textit{Retrieval-Augmented Generation}) d'aide à
l'enquête judiciaire déployé au CFIA, Gendarmerie nationale. Il traite des
documents hétérogènes (images, audio, vidéo, textes) pour permettre aux
enquêteurs d'interroger leur corpus documentaire en langage naturel.

En vertu de l'AI Act (Règlement UE 2024/1689), IDeXtend sera probablement classifié
\textbf{système à haut risque} au titre de l'Annexe III §6 (application de la
loi). Ce dossier documente scientifiquement le \textbf{système de scoring de
confiance} de chaque brique du pipeline, avec pour chaque composant : la
formule mathématique du score, sa justification par l'état de l'art.

\bigskip
\noindent\textbf{État d'implémentation (mai 2026) :}

\begin{center}
\begin{tabular}{llll}
\toprule
\textbf{Composant} & \textbf{Modèle} & \textbf{Scoring implémenté} & \textbf{Validé} \\
\midrule
Détection langue & FastText LID            & \textbf{Oui}  & Non \\
Traduction       & NLLB-600M               & \textbf{Oui}  & Non \\
Reranker         & BAAI/bge-reranker-v2-m3 & \textbf{Oui}  & Non \\
OCR              & QwenVL 2.5-7B           & À définir     & Non \\
ASR (Whisper)    & whisper-large-v3        & À implémenter & Non \\
LLM (Qwen3)      & Qwen3-4B                & À implémenter & Non \\
Embeddings       & BAAI/bge-m3             & N/A           & Non \\
\bottomrule
\end{tabular}
\end{center}

\clearpage

% --- Composants -----------------------------------------------------------
\chapter{Détection automatique de langue}

\section{Identification}

\begin{center}
\begin{tabular}{ll}
\toprule
Modèle         & \texttt{facebook/fasttext-language-identification} \\
Nb langues     & 217 \\
Statut scoring & \textbf{Implémenté} \\
Fichier        & \texttt{DataPipelines/LanguageDetection.py} \\
\bottomrule
\end{tabular}
\end{center}

\section{Fonctionnalité}

Premier composant de la chaîne : détermine si une traduction est nécessaire.
Déclenché sur les 700 premiers + 1500--2500 caractères du texte. Retourne le
code langue BCP-47 FLORES-200 (ex.\ \texttt{fra\_Latn}, \texttt{ara\_Arab}).

\section{Fonction de scoring de confiance}

\noindent\textit{Implémentation dans
\texttt{LanguageDetection.compute\_confidence()}.}

\medskip
FastText applique un softmax sur 217 classes. Le score de confiance est
directement la probabilité de la langue prédite :

\begin{equation}
  s_{\text{LID}}(x)
  = P(\hat{l} \mid x)
  = \max_{l \in \mathcal{L}} P(l \mid x)
  \;\in (0, 1]
  \label{eq:lid}
\end{equation}

Les top-2 prédictions sont stockées dans les métadonnées pour auditabilité.

\paragraph{Seuils.}

\begin{center}
\begin{tabular}{lll}
\toprule
$s \geq 0.90$ & Langue certaine & Décision automatique fiable \\
$0.60 \leq s < 0.90$ & Langue probable & Logging, surveiller \\
$s < 0.60$ & Détection ambiguë & Ne pas traduire automatiquement \\
\bottomrule
\end{tabular}
\end{center}

\textbf{Métadonnées loguées :} \texttt{detected\_lang},
\texttt{top\_k} (liste des 2 meilleures prédictions avec scores).

\section{Justification scientifique}

FastText, probabilités softmax~\cite{joulin2017bag} ;
cas difficiles (code-switching, dialectes)~\cite{caswell2020language}.

\section{Validation empirique}

Aucun dataset d'évaluation annoté n'est disponible à ce jour. Le
\texttt{compute\_confidence()} est actuellement validé par inspection manuelle
sur des cas réels et des scénarios limites (\textit{smoke testing}).

\noindent\textbf{Métriques à mesurer lors de la constitution du dataset :}
\begin{itemize}
  \item Corrélation entre le score de confiance et la bonne classification ;
  \item Calibration (ECE) : un score de confiance de 0.8 correspond-il à
    environ 80\,\% de réussite sur le dataset ?
  \item Pouvoir discriminant (AUC-ROC) : le score sépare-t-il les prédictions
    correctes des incorrectes ?
\end{itemize}

\noindent\textbf{Statut :} \textit{À implémenter.}

\paragraph{Stress tests à conduire.}
Textes bilingues mélangés, textes très courts ($< 20$ caractères), arabe
dialectal levantin, translittération arabe en latin.

\chapter{Traduction automatique --- NLLB-600M}

\section{Identification}

\begin{center}
\begin{tabular}{ll}
\toprule
Modèle         & \texttt{facebook/nllb-200-distilled-600M} \\
Beam search    & $B = 5$ faisceaux \\
Statut scoring & \textbf{Implémenté} \\
Fichier        & \texttt{DataPipelines/TranslationModel.py} \\
\bottomrule
\end{tabular}
\end{center}

\section{Benchmark FLORES+ (200 phrases, 8 paires de langues)}

\begin{center}
\begin{tabular}{lrrrrr}
\toprule
\textbf{Modèle} & FR$\to$EN & EN$\to$FR & FR$\to$AR & AR$\to$FR & ms/phrase \\
\midrule
\textbf{NLLB-600M} & 41.0 & \textbf{46.3} & 13.9 & \textbf{29.7} & \textbf{130--180} \\
Qwen3.5-4B & \textbf{41.2} & 44.7 & 14.5 & 29.3 & 660--900 \\
TGemma-4B  & 37.9 & 40.4 & \textbf{14.5} & 26.4 & 620--940 \\
\bottomrule
\end{tabular}
\end{center}

NLLB retenu : gagne 5--6/8 directions, 4--5$\times$ plus rapide.
Résultats complets dans \texttt{scoring\_rayan/\_archive/}.

\section{Fonction de scoring de confiance}

\noindent\textit{Implémentation dans \texttt{TranslationModel.compute\_confidence()}.}

\medskip
Soit $K$ le nombre de chunks, $\ell_k$ le log-prob normalisé best-beam,
$p_k = \exp(\ell_k)$, et $r_k = \text{output\_len}_k / \text{input\_len}_k$
le ratio de longueur.
%definir P(barre)

\begin{equation}
  \phi = 1 - \frac{|\{k : r_k < 0.3 \;\text{ou}\; r_k > 4\}|}{K} \cdot 0.5
  \label{eq:trad_penalty}
\end{equation}

\begin{equation}
  s_{\text{trad}}(x, \hat{y})
  = \bigl(0.7 \cdot \bar{p} + 0.3 \cdot \min_k p_k\bigr) \cdot \phi
  \;\in (0, 1]
  \label{eq:trad}
\end{equation}

\paragraph{Seuils.}

\begin{center}
\begin{tabular}{lll}
\toprule
$s \geq 0.70$ & Traduction fiable & Exploitable \\
$0.40 \leq s < 0.70$ & Traduction incertaine & Vérification recommandée \\
$s < 0.40$ & Très incertaine & Relecture humaine \\
\bottomrule
\end{tabular}
\end{center}

\textbf{Métadonnées loguées :} \texttt{num\_chunks}, \texttt{src\_lang},
\texttt{tgt\_lang}, \texttt{mean\_beam\_prob}, \texttt{min\_beam\_prob},
\texttt{outlier\_chunks}.

\section{Justification scientifique}

NLLB~\cite{nllb2022} ; calibration beam NMT~\cite{ott2018analyzing} ;
Quality Estimation non supervisée~\cite{fomicheva2020unsupervised}.

\section{Validation empirique}

Aucun dataset d'évaluation annoté n'est disponible à ce jour. Le
\texttt{compute\_confidence()} est actuellement validé par inspection manuelle
sur des cas réels et des scénarios limites (\textit{smoke testing}).

\noindent\textbf{Métriques à mesurer lors de la constitution du dataset :}
\begin{itemize}
  \item Corrélation entre le score de confiance et la métrique de qualité
    appropriée (ex. BLEU, chrF++) ;
  \item Calibration (ECE) : un score de confiance de 0.8 correspond-il à
    environ 80\,\% de réussite sur le dataset ?
  \item Pouvoir discriminant (AUC-ROC) : le score sépare-t-il les traductions
    correctes des incorrectes ?
\end{itemize}

\noindent\textbf{Statut :} \textit{À implémenter.}


% # Composite: mean weighted 0.7, min weighted 0.3 (penalises worst-case chunks)
%        score = (0.7 * mean_beam_prob + 0.3 * min_beam_prob) * outlier_penalty
% la formule de calcul est aussi composite a étudié
\chapter{Reranker}

\section{Identification}

\begin{center}
\begin{tabular}{ll}
\toprule
Modèle         & \texttt{BAAI/bge-reranker-v2-m3} (CrossEncoder) \\
Statut scoring & \textbf{Implémenté} \\
Fichier        & \texttt{DataPipelines/Reranker.py} \\
\bottomrule
\end{tabular}
\end{center}

\section{Fonctionnalité}

Le reranker reçoit les 25 chunks candidats des embeddings et les reclasse par
pertinence à la question en traitant chaque paire (question, chunk) ensemble.
Il filtre les chunks sous le seuil \texttt{minimum\_acceptable\_score = 0.05}
selon la méthode de sélection configurée (\texttt{smart}, \texttt{kmeans},
\texttt{top}, \texttt{top\_kmeans}).

\section{Fonction de scoring de confiance}

\noindent\textit{Implémentation dans \texttt{DataPipelines/Reranker.py}.}

\medskip
Le cross-encoder produit un logit $l_i$ par paire (question, chunk), converti
en score $\tilde{s}_i = \sigma(l_i) \in (0,1)$.

Le score de confiance actuel est une moyenne pondérée entre la moyenne des
chunks retenus et le plus faible d'entre eux :

\begin{equation}
  s_{\text{reranker}}
  = 0.7 \cdot \bar{s} + 0.3 \cdot \min_{i \in \mathcal{I}} \tilde{s}_i
  \label{eq:reranker}
\end{equation}

\noindent\textbf{Cette formule est temporaire.} La pondération 0.7/0.3 n'a
pas été validée empiriquement. Elle ne tient pas compte de la taille du corpus :
avec peu de documents disponibles, le système peut retenir des chunks
faiblement pertinents faute de mieux, ce que le score ne reflète pas.

\textbf{Cas particuliers.}
\begin{itemize}
  \item $|\mathcal{I}| = 0$ : $s = 0.0$ (aucun candidat retenu).
  \item $|\mathcal{I}| = 1$, fast path : $s = 0.5$ (non mesuré par convention).
\end{itemize}

\paragraph{Seuils.}

\begin{center}
\begin{tabular}{lll}
\toprule
$s \geq 0.70$ & Bons candidats & Réponse LLM probablement pertinente \\
$0.40 \leq s < 0.70$ & Pertinence moyenne & Vérification recommandée \\
$s < 0.40$ & Peu pertinent & Avertissement \\
\bottomrule
\end{tabular}
\end{center}

\noindent Ces seuils sont provisoires et non calibrés.

\textbf{Métadonnées loguées :} \texttt{num\_selected}, \texttt{mean\_score},
\texttt{min\_score}, \texttt{max\_score}, \texttt{per\_index\_scores}.


\noindent\textit{Tests de régression automatisés :
\texttt{Tests/test\_reranker\_regression.py}.}

\chapter{OCR --- Reconnaissance Optique de Caractères}

\section{Identification}

\begin{center}
\begin{tabular}{ll}
\toprule
\textbf{Champ} & \textbf{Valeur} \\
\midrule
Modèle actif    & QwenVL 2.5-7B (vLLM) \\
Statut scoring  & À implémenter \\
Fichier         & \texttt{DataPipelines/VLMModel.py} \\
\bottomrule
\end{tabular}
\end{center}

\section{Fonctionnalité}

Le composant OCR convertit les images en texte exploitable. Il est déclenché
sur chaque n\oe ud image portant le flag \texttt{OCR\_FLAG}. Le texte produit
est découpé en chunks, vectorisé par BGE-M3, traduit si nécessaire, puis
indexé pour la recherche.

\textbf{Entrée.} Image binaire (JPEG, PNG, BMP) $> 100 \times 100$ px, ratio
max/min $< 10$.\\
\textbf{Sortie.} Texte brut multi-lignes, ordre de lecture reconstitué.

\section{Modèles évalués}

\begin{center}
\begin{tabular}{llll}
\toprule
\textbf{Modèle} & \textbf{Version} & \textbf{Benchmarké} & \textbf{Retenu} \\
\midrule
EasyOCR + Tesseract & 1.7 + 5.x & Non (legacy)   & Non \\
QwenVL 2.5-7B       & 7B Q8     & Non            & Oui \\
Marker              & TDB       & Non        & Non \\
Docling (IBM)       & TDB       & Non        & Non \\
InternVL 3.5-7B     & TDB       & Non        & Non \\
\bottomrule
\end{tabular}
\end{center}

\section{Protocole d'évaluation}

\noindent\textbf{Métriques.} TBD.

\noindent\textbf{Dataset à constituer.}  TBD 

\section{Fonction de scoring de confiance}

\subsection{Mode legacy --- EasyOCR + Tesseract}

%Soit $W^{+} = \{w_i : c_i \geq 0\}$ les mots détectés avec score Tesseract
%non nul.

% \begin{equation}
%   s_{\text{legacy}}(x, \hat{y})
%   = \frac{1}{|W^{+}|} \sum_{i \in W^{+}} \frac{c_i}{100}
%   \;\in [0, 1]
%   \label{eq:ocr_legacy}
% \end{equation}

% La rotation de l'image retient la transcription maximisant ce score.
TBD
\subsection{Mode VLM / QwenVL (actif) --- Proposition}

% Le VLM ne produit pas de scores par token via l'API actuelle. Deux approches
% sont proposées.

% \paragraph{Option A --- Confiance verbalisée (recommandée).}
% Ajout d'une instruction de post-génération~:
% \textit{\og Transcris le texte. Termine par : CONFIANCE: <entier 0--100>.\fg{}}

% \begin{equation}
%   s_{\text{VLM}}(x, \hat{y})
%   = \frac{c_{\text{verb}}}{100} \cdot \delta_{\text{format}}
%   \;\in [0, 1]
%   \label{eq:ocr_vlm}
% \end{equation}

% où $\delta_{\text{format}} = 1$ si le token CONFIANCE est présent et
% parseable, $0.5$ sinon.

% \paragraph{Option B --- Self-Consistency ($K$ inférences).}

% \begin{equation}
%   s_{\text{SC}}(x)
%   = 1 - \frac{1}{K(K-1)} \sum_{i \neq j} \mathrm{NED}(\hat{y}_i, \hat{y}_j)
%   \;\in [0, 1]
%   \label{eq:ocr_sc}
% \end{equation}

% Coût : $K \times$ latence.
TBD
\paragraph{Seuils provisoires.}

\begin{center}
\begin{tabular}{lll}
\toprule
\textbf{Plage} & \textbf{Interprétation} & \textbf{Action} \\
\midrule
TBD           & Confiance moyenne & Relecture recommandée \\
TBD           & Haute confiance  & Exploitable \\
TBD           & Faible confiance  & Relecture obligatoire \\
\bottomrule
\end{tabular}
\end{center}

\section{Justification scientifique}

\section{Validation empirique}


\chapter{ASR --- Reconnaissance Automatique de la Parole}

\section{Identification}

\begin{center}
\begin{tabular}{ll}
\toprule
\textbf{Champ} & \textbf{Valeur} \\
\midrule
Modèle         & \texttt{openai/whisper-large-v3} \\
Mode           & Chunked (30\,s/chunk, batch = 8) \\
Statut scoring & À implémenter \\
Fichier        & \texttt{DataPipelines/WhisperModel.py} \\
\bottomrule
\end{tabular}
\end{center}

\section{Fonctionnalité}

Whisper transcrit les fichiers audio et vidéo en texte. Dans le contexte
judiciaire : écoutes téléphoniques, auditions, vidéos de surveillance.

\section{Modèles évalués}

\begin{center}
\begin{tabular}{llll}
\toprule
\textbf{Modèle} & \textbf{Version} & \textbf{Benchmarké} & \textbf{Retenu} \\
\midrule
whisper-large-v3   & large-v3      & Non  & \textbf{Oui} \\
whisper-medium     & medium        & Non  & Non \\
Faster-Whisper CT2 & CT2 large-v3  & Non  & Non \\
\bottomrule
\end{tabular}
\end{center}

\section{Fonction de scoring de confiance}

\texttt{WhisperModel.infer()} retourne uniquement \texttt{result["text"]} ---
aucun beam score ni métadonnée de confiance n'est exposé dans l'état actuel.

\textbf{Pour implémenter un scoring}, il faudrait modifier \texttt{infer()} afin
de récupérer les \texttt{output\_scores} de la pipeline HuggingFace.
Une approche cohérente avec \texttt{TranslationModel} serait d'utiliser
les log-probs best-beam par chunk (même architecture seq2seq, même beam search).
Cette piste reste à valider --- aucune formule n'est fixée à ce stade.

\section{Justification scientifique}

Architecture Whisper~\cite{radford2022whisper}.

\section{Validation empirique}

Aucun dataset d'évaluation annoté n'est disponible à ce jour. Le
\texttt{compute\_confidence()} est actuellement validé par inspection manuelle
sur des cas réels et des scénarios limites (\textit{smoke testing}).

\noindent\textbf{Métriques à mesurer lors de la constitution du dataset :}
\begin{itemize}
  \item Corrélation entre le score de confiance et la métrique de qualité
    appropriée (ici $1 - \mathrm{WER}$) ;
  \item Calibration (ECE) : un score de confiance de 0.8 correspond-il à
    environ 80\,\% de réussite sur le dataset ?
  \item Pouvoir discriminant (AUC-ROC) : le score sépare-t-il les transcriptions
    correctes des incorrectes ?
\end{itemize}

\noindent\textbf{Statut :} \textit{À implémenter.}

\chapter{LLM --- Grand Modèle de Langage (Qwen3)}

\section{Identification}

\begin{center}
\begin{tabular}{ll}
\toprule
Config active  & \texttt{Qwen3-4B-Instruct-ctx-8k-30-user} \\
Modèle         & \texttt{Qwen/Qwen3-4B} \\
Backend        & vLLM \\
Contexte       & 8\,000 tokens \\
Statut scoring & \textbf{À implémenter} \\
Fichier        & \texttt{DataPipelines/LlmModel.py} \\
\bottomrule
\end{tabular}
\end{center}

\section{Fonctionnalité}

Le LLM est le composant de génération finale : réponses aux enquêteurs,
résumés, extraction d'entités, graphe de connaissances.

\section{Pourquoi les log-probs seuls sont insuffisants}

Pour un LLM instruit par RLHF/DPO, la probabilité de séquence
$P(\hat{y}|x) = \prod_t P(y_t | y_{<t}, x)$ est :
\begin{enumerate}
  \item \textbf{Mal calibrée} post-RLHF (scores absolus peu interprétables).
  \item \textbf{Biaisée par la longueur} : les réponses courtes ont
        mécaniquement des log-probs plus élevés.
  \item \textbf{Non discriminante des hallucinations} : un modèle peut être
        très confiant en générant une information fausse.
\end{enumerate}

Le score de confiance du LLM ne doit pas être interprété comme une
probabilité de vérité. Il est utilisé comme un signal de triage pour détecter
les réponses nécessitant une relecture humaine.

\section{Approche de scoring proposée}

\paragraph{Signal 1 --- \textit{Retrieval quality}.}

On mesure si les chunks fournis au LLM contiennent les éléments nécessaires
pour répondre à la question.

\begin{equation}
  s_{\text{retrieval}}(q, C)
  = \mathbb{1}[\text{gold chunk} \in C_{1:k}]
  \label{eq:llm_retrieval}
\end{equation}

En absence de dataset annoté, ce signal est estimé par inspection humaine sur
un échantillon de questions.

\paragraph{Signal 2 --- \textit{Groundedness}.}

La réponse est découpée en affirmations atomiques. Chaque affirmation est
classée comme \textit{supported}, \textit{unsupported} ou \textit{contradicted}
par les chunks fournis.

\begin{equation}
  s_{\text{groundedness}}(\hat{y}, C)
  = \frac{\#\text{\textit{supported claims}}}{\#\text{\textit{total claims}}}
  \label{eq:llm_groundedness}
\end{equation}

\begin{equation}
  r_{\text{contradiction}}(\hat{y}, C)
  = \frac{\#\text{\textit{contradicted claims}}}{\#\text{\textit{total claims}}}
  \label{eq:llm_contradiction}
\end{equation}

\paragraph{Signal 3 --- \textit{Abstention}.}

Le système doit refuser de conclure lorsque les documents ne contiennent pas
l'information demandée.

\begin{equation}
  s_{\text{abstain}}
  = \mathbb{1}[
    \text{réponse absente des documents}
    \Rightarrow
    \text{\textit{abstention} ou réponse qualifiée}
  ]
  \label{eq:llm_abstain}
\end{equation}

\paragraph{Score de triage.}

\begin{equation}
  s_{\text{triage}}
  = \min\!\Bigl(1,\;
    \alpha\, s_{\text{groundedness}}
    + \beta\, s_{\text{retrieval}}
    - \lambda\, r_{\text{contradiction}}
  \Bigr) \cdot s_{\text{abstain}}
  \label{eq:llm_triage}
\end{equation}

avec $\alpha + \beta = 1$, $\lambda \geq 0$.
Le facteur $s_{\text{abstain}} \in \{0, 1\}$ agit comme un veto :
si le système devait s'abstenir et ne l'a pas fait, le score est nul.
Les coefficients $\alpha, \beta, \lambda$ ne sont pas fixés a priori ;
ils doivent être calibrés sur le dataset de validation~\cite{guo2017calibration}.

\paragraph{Règle opérationnelle.}

\begin{center}
\begin{tabular}{lll}
\toprule
\textbf{Condition} & \textbf{Interprétation} & \textbf{Action} \\
\midrule
Preuve absente                        & Réponse \textit{ungrounded}  & \textit{Abstention} obligatoire \\
Contradiction détectée                & Risque élevé                 & Relecture humaine \\
$s_{\text{groundedness}} < 0.8$       & \textit{Grounding} insuffisant & Relecture humaine \\
Toutes les \textit{claims} supportées & Réponse exploitable          & Citation obligatoire \\
\bottomrule
\end{tabular}
\end{center}

\section{Validation empirique}

\subsection*{Dataset cible}

\begin{center}
\begin{tabular}{ll}
\toprule
Affaires            & $\sim$10 cas réels anonymisés \\
Documents           & $\sim$50 documents au total \\
Questions annotées  & schéma ci-dessous \\
\bottomrule
\end{tabular}
\end{center}

\begin{verbatim}
{
  "case_id":          "affaire_001",
  "question":         "Quel véhicule était impliqué ?",
  "reference_answer": "Un Renault Clio immatriculé AB-123-CD",
  "golden_chunks":    ["chunk_002", "chunk_006"],
  "answerable":       true,
  "ambiguous":        false
}
\end{verbatim}

Le dataset doit être constitué et verrouillé avant tout réglage des seuils
et coefficients.

\noindent\textbf{Champs calculés à l'évaluation} (non dans le dataset statique) :
\texttt{retrieved\_chunks}, \texttt{generated\_answer}, \texttt{abstained},
\texttt{human\_score} (jugement humain 0--1, proxy pour $s_{\text{groundedness}}$
en l'absence d'annotation \textit{claim}-level).

\subsection*{Métriques}

\begin{itemize}
  \item \textbf{Retrieval hit@k} : le chunk annoté pertinent est-il dans
        les $k$ chunks fournis au LLM ?
  \item \textbf{Exactitude de réponse} : réponse correcte, partielle,
        incorrecte ou abstention (jugement humain).
  \item \textbf{\textit{Unsupported claims rate}} : proportion d'affirmations
        absentes des sources.
  \item \textbf{\textit{Contradiction rate}} : proportion d'affirmations
        contredites par les sources.
  \item \textbf{\textit{Abstention} correcte} : le système refuse-t-il de
        conclure sur les questions négatives et ambiguës ?
\end{itemize}

\noindent\textbf{Statut actuel :} validation manuelle exploratoire uniquement.
Le système ne doit pas être présenté comme calibré tant que ce dataset
n'a pas été constitué.


% Pour les logs probs, on prend des articles existant, et on va chercher des infomrations precises, chercher des elements precis, quand on sait que la reponse est la, quelle est le comportement de logs probs, quand la reposne nest pas la 
% qunad la question est mla posé, qu'est-ce qui se passe, 
% on fait une vingtaine de questions et le but s'est de voir comment se comporte le logsprobs,
% la quesiton n'est pas adapté au texte, et les logs probs sont élevé -> c'est pas bon 
% prendre des articles de journaux Emmanuel, et poser des questions sur les faits, les personnes, les lieux, et voir comment se comporte les logs probs, est ce que le score de confiance est corrélé avec la qualité de la réponse
% Soit infomration precises, soit large, question pas adaptée vs question adaptée 



\chapter{Embeddings --- Représentation vectorielle}

\section{Identification}

\begin{center}
\begin{tabular}{ll}
\toprule
Modèle         & \texttt{BAAI/bge-m3} \\
Dimension      & 1\,024 \\
Statut scoring & À implémenter \\
Fichier        & \texttt{DataPipelines/EmbeddingsModel.py} \\
\bottomrule
\end{tabular}
\end{center}

\section{Fonctionnalité}

Le modèle d'embeddings encode chunks et questions en vecteurs denses.
La documentation du benchmark de sélection est dans
\texttt{documentation\_pierre/choosing\_embeddings.md}.

\section{Benchmark (extrait)}

\begin{center}
\begin{tabular}{lll}
\toprule
\textbf{Modèle} & \textbf{Position moy. réponse} & \textbf{Retenu} \\
\midrule
paraphrase-multilingual-MiniLM-L12-v2 & 31    & Non \\
\textbf{BAAI/bge-m3}                  & \textbf{4.7} & \textbf{Oui} \\
\bottomrule
\end{tabular}
\end{center}

\noindent Score = rang moyen de la réponse correcte (plus bas = mieux).
Protocole \og needle in the haystack\fg{}, 4 langues, script
\texttt{evaluate\_embeddings.py}.

\section{Scoring de confiance}

\texttt{EmbeddingsModel} est un encodeur pur : il produit un vecteur,
sans score de confiance associé. Le retrieval est effectué en aval par
\texttt{CypherHandler.search\_embeddings()} dans \texttt{LargeLanguageHandler/rag.py}.

Un \texttt{compute\_confidence()} au niveau de l'encodeur n'a pas de sens.
La qualité du retrieval est évaluée via le benchmark existant (position moyenne
de la réponse correcte, script \texttt{evaluate\_embeddings.py}) et via la
métrique \textbf{retrieval hit@k} du dataset d'évaluation LLM (chapitre~5).


% Daniel : variation de phrase, utilisation de synonyme, et avoir un score pour l'embedding 
% Comparaison d'embeddings : pouruqoi est-ce qu'on a choisi celui-ci ? 
% L'embeddings est un sous produit, l'utilsateur ne le voit pas directement, c'est un composant intermédiaire.
% il n'a pas besoin de savoir 


% Par contre embedding d'image si,
% -> score de match entre la query est l'image, 

% --- Synthèse -------------------------------------------------------------
\chapter*{Synthèse et prochaines étapes}
\addcontentsline{toc}{chapter}{Synthèse et prochaines étapes}

Trois composants ont déjà leur \texttt{compute\_confidence()} implémentée :
Reranker, Traduction (NLLB) et Détection de langue (FastText). Les quatre
composants restants nécessitent une implémentation, dont le LLM qui requiert
une approche composite (ancrage contextuel + confiance verbalisée).

\paragraph{Actions prioritaires.}
\begin{enumerate}
  \item Constituer les datasets d'évaluation annotés (section 4 de chaque fiche).
  \item Implémenter \texttt{compute\_confidence()} pour OCR, ASR, Embeddings et LLM.
  \item Conduire les expériences de validation empirique : Spearman $\rho$, ECE, AUC-ROC.
  \item Compléter le benchmark OCR : Marker et Docling vs QwenVL 2.5-7B.
\end{enumerate}

\paragraph{Limite structurelle : explicabilité et résistance au re-benchmark.}

IDeXtend est conçu pour être modulaire : chaque composant est remplaçable
indépendamment. Cette modularité est un atout technique, mais elle crée une
tension avec les exigences d'explicabilité de l'AI Act.

Valider un composant nécessite un dataset annoté, du temps d'annotation, et
une infrastructure d'évaluation. En pratique, lorsqu'un utilisateur ou un
opérateur souhaite substituer un modèle par une version plus récente --- même
si le remplacement est techniquement transparent (\textit{plug and play}) ---
la résistance métier et terrain à relancer un benchmark complet est réelle.
Les raisons sont multiples : coût d'annotation, indisponibilité des données
de test originales, pression opérationnelle, sentiment que ``le nouveau modèle
est forcément meilleur''.

Ce risque est structurel : un système peut dériver silencieusement dès lors
qu'un composant est mis à jour sans re-validation. L'AI Act (Art.~9 et 15)
exige pourtant une réévaluation à chaque changement substantiel.

\textbf{Piste.} Définir contractuellement, pour chaque composant, un
\textit{smoke test} minimal automatisé (quelques cas annotés en dur dans
le dépôt) qui tourne à chaque mise à jour de modèle en CI/CD. Ce n'est pas
une validation complète, mais c'est une barrière minimale contre les
régressions silencieuses et un signal d'alerte si le comportement change.

\bibliographystyle{plain}
\bibliography{references}

\end{document}
